% SmallSat Europe + JOSS Mega Paper Draft
% 3-Magnetorquer, 1-Reaction-Wheel Architecture Demonstration

% ============================================================
% TODO SUMMARY BY VENUE
% ============================================================
%
% FOR SMALLSAT EUROPE (practical, accessible, ~6 pages):
%   - KEEP: Introduction with CONOPS motivation
%   - KEEP: Mission applications section (1U, constellation, eclipse)
%   - KEEP: Monte Carlo results tables (summary only)
%   - KEEP: BeaverCube-2 real mission example
%   - TRIM: Detailed control allocation math (reference other paper)
%   - TRIM: Extensive prior work review
%   - TRIM: Momentum management details (brief mention only)
%   - ADD: "When to use 3+1" decision flowchart
%   - ADD: SWaP-C comparison table prominently
%   - ADD: Practitioner-focused recommendations
%
% FOR JOSS (software/methodology focus, ~8-10 pages):
%   - KEEP: All technical content
%   - KEEP: Monte Carlo methodology details
%   - KEEP: Control approach section with allocation math
%   - ADD: Software availability statement
%   - ADD: Code snippets showing configuration
%   - ADD: Reproducibility information (random seeds, parameters)
%   - ADD: Comparison with Basilisk implementation effort
%   - TRIM: Mission cost analysis (less relevant for software paper)
%
% ============================================================
% *** CRITICAL DATA GAP ASSESSMENT (2026-01-24 Session 4) ***
% ============================================================
%
% PROBLEM: Current figures in experiment_outputs/ are NOT publication-ready!
%
% What experiment_outputs/3p1_paper/ actually contains:
%   - Duration: 200 seconds (TOO SHORT - MTQs need 500-1000s to converge)
%   - Controller: PD only (NO ALTRO planner)
%   - Trials: 10 (NOT statistically significant)
%   - Results: 3+0=87°, 3+1=63°, 3+3=7° (WRONG - too short to converge)
%
% What papers/3MTQ+1RW/output_data/ contains (ALTRO + TVLQR, 10 trials, 500s):
%   Mean: 30°, range 3.6°-78° (DOES NOT MATCH THESIS)
%
% What thesis claims (planning.tex):
%   - 3MTQ+1RW: 96% within 1° for single slew
%   - Mean final error: 0.45° (median 0.03°)
%   THIS IS ~60x BETTER THAN CURRENT DATA
%
% REQUIRED ACTION: Re-run experiments with proper settings:
%   - Duration: 1000s minimum
%   - Trials: 100 minimum
%   - Goal type: REDUCED-ATTITUDE (thesis used this!)
%   - ALTRO params from ALTRO_TUNING_NOTES.md Session 6
%
% ============================================================
% DATA GENERATION TODO LIST
% ============================================================
% ** CRITICAL - BEFORE SUBMISSION **
% TODO-DATA-0: REPRODUCE THESIS RESULTS with proper parameters
%
% ** HIGH PRIORITY **
% TODO-DATA-1: Verify Monte Carlo numbers match simulation data
%    - TEST STRUCTURE:
%      * First comparison: 3+0, 3+1, 3+3 ALL with PD-based control
%      * Then: 3+1 with planner to show improvement over PD
%      * Optional: 3+3 with planner (would also do well, consider including)
%    - Numbers to verify:
%      * 2.3 deg mean (3+1 PD), 0.05 deg (3+1 planner)
%      * 73% within 1 deg (PD), 100% (planner)
%    - NOTE: 3+3 with planner would outperform 3+1 with planner, but
%      the key comparison is 3+1+planner vs 3+3+PD (shows 3+1 can match!)
%
%    *** THESIS MONTE CARLO RESULTS (planning.tex) - USE THESE ***
%    Single 180° slew results:
%      * MTQ-only: 73% within 10° for single slew
%      * 3MTQ+1RW: 96% within 1° for single slew
%    Goal-set (reduced-attitude) slews:
%      * MTQ-only: 67% within 1°
%      * 3MTQ+1RW: 96% within 1°
%    Multi-target (3 targets, 500s simulation):
%      * 3MTQ+1RW: 98%+ within 10° for each target
%      * Mean final attitude error: 0.45° (median 0.03°)
%    Sequential trajectory planning (thesis Section 7.4):
%      * Sub-degree pointing accuracy demonstrated
%      * 96% convergence rate across varied initial conditions
%
% DATA SOURCES (from Generalized_ADCS codebase):
%    - Monte Carlo infrastructure: testing/paper_todo_tests/test_todo_sim_monte_carlo.py
%    - Allocation comparison results: research/allocation_comparison_*.json
%    - Closed-loop results: research/closed_loop_*.json
%    - LP vs QP results (from codebase):
%      * LP: 17.02° final error, 0.0036° direction error, ~26% magnitude
%      * QP: 25.70° final error, 33.01° direction error, ~60% magnitude
%    - Spacecraft configs: testing/paper_todo_tests/test_todo_data_generation.py
%      * 3U_Standard: mass=4.0kg, inertia=diag([0.035, 0.035, 0.007]) kg⋅m²
%
% TODO-DATA-2: Generate momentum management comparison [NOT YET AVAILABLE]
%    - Continuous vs scheduled desaturation
%    - Wheel momentum evolution over multiple orbits
%    - Pointing accuracy during desaturation windows
%    - This data needs to be generated
%
% TODO-DATA-3: Generate goal formulation comparison data
%    - Full-attitude vs reduced-attitude for 3+1
%    - Quantify improvement percentage
%
% TODO-BC2-1: Get BeaverCube-2 mission parameters
%    - Mass, inertia
%    - Actuator specs (MTQ dipole limits, RW specs)
%    - Pointing requirements
%    - What constraint drove 3+1 choice? (mass? power? volume? cost?)
%
% TODO-BC2-2: Get BeaverCube-2 expected performance
%    - Simulation results if available
%    - Expected pointing accuracy
%    - Comparison to what would have been achieved with 3+0 or 3+3
%
% TODO-BC2-3: Document BeaverCube-2 CONOPS
%    - Imaging targets and timing
%    - Slew requirements
%    - Eclipse handling
%
% ** MEDIUM PRIORITY **
% TODO-DATA-4: Run actuator failure simulations
%    - 3+1 with wheel failure (reverts to 3+0)
%    - Compare graceful degradation performance
%
% TODO-DATA-5: Generate sensitivity analysis
%    - Inertia errors, magnetic field errors
%    - Actuator scaling errors
%
% ** LOW PRIORITY **
% TODO-DATA-6: Generate spinning solution examples
%    - Cases where planner discovers spin maneuvers
%    - Disturbance rejection via controlled spin
%
% ============================================================
% FIGURE TODO LIST
% ============================================================
% TODO-FIG-1: Actuator configuration comparison (3+0, 3+1, 3+3) - visual
% TODO-FIG-2: Monte Carlo pointing error distributions (histogram/CDF)
% TODO-FIG-3: Time series showing 3+1 vs 3+0 vs 3+3 convergence
% TODO-FIG-4: Torque envelope comparison for each configuration
% TODO-FIG-5: Momentum evolution with continuous vs scheduled desat
% TODO-FIG-6: Decision flowchart for architecture selection
%
% ============================================================
% LITERATURE REVIEW TODO LIST (from thesis background.tex)
% ============================================================
%
% ** HYBRID ACTUATOR ARCHITECTURES **
% TODO-LIT-1: Document that 3+1 is a "novel actuator combination" rarely studied
%             - Thesis calls this out explicitly
%             - Most literature is 3+0 (pure magnetic) or 3+3 (full RW)
%             - Little work on hybrid underactuated configurations
%
% ** MAGNETIC CONTROL FOUNDATIONS **
% TODO-LIT-2: Ovchinnikov & Roldugin "Survey of Magnetic Control" (2019)
%             - Comprehensive survey, positions 3+1 relative to prior art
% TODO-LIT-3: Lovera & Astolfi (2004, 2005, 2006)
%             - Projection-based magnetic control, stability proofs
% TODO-LIT-4: Wisniewski - Ørsted sliding mode (1997, 1999)
%             - First systematic magnetic control treatment
%
% ** FLIGHT-PROVEN HYBRID SYSTEMS **
% TODO-LIT-5: PRISMA TANGO (Chasset 2006, 2013)
%             - Flight-proven magnetic control, hybrid with momentum wheel
% TODO-LIT-6: GOCE (Sechi 2005)
%             - 3-axis magnetic control, high accuracy
%
% ** TRAJECTORY PLANNING **
% TODO-LIT-7: ALTRO (Howell 2019, Jackson 2021)
%             - Trajectory optimization algorithm used for planner
% TODO-LIT-8: Gatherer "Magnetorquer-Only ALTRO" (2019)
%             - Prior ALTRO application to magnetic control
%
% ** UNDERACTUATED CONTROL **
% TODO-LIT-9: Petersen & Kolmanovsky - 2RW MPC (2015, 2017)
%             - Underactuated reaction wheel control
% TODO-LIT-10: Inamori - Two-step rotation (2016)
%              - Alternative underactuated approach
%
% ** SMALL SATELLITE CONTEXT **
% TODO-LIT-11: Cappelletti "CubeSat Handbook" (2020)
%              - CubeSat ADCS constraints, actuator options
% TODO-LIT-12: Cite RW failure statistics (~50% of ADCS failures)
%              - Justifies reducing RW count
%
% ============================================================
% KEY CITATIONS FOR 3+1 PAPER TEXT:
% ============================================================
% Introduction:
%   - "Magnetic control has advantages [Ovchinnikov2019] but is underactuated"
%   - "Most hybrid approaches use 3+3 [Markley2014], we propose 3+1"
% Background:
%   - Magnetic control theory [Lovera2004], [Wisniewski1997]
%   - Flight heritage [Chasset2013], [Sechi2005]
% Planner:
%   - ALTRO algorithm [Howell2019], [Jackson2021]
%   - Prior MTQ application [Gatherer2019]
% ============================================================
%
% ============================================================
% DETAILED EXPERIMENT LIST (Priority Order)
% ============================================================
%
% === EXPERIMENT SET A: Core Architecture Comparison (ESSENTIAL) ===
% Run identical conditions across 3+0, 3+1, 3+3 to produce Tables 1-2
%
% PRIMARY SPACECRAFT: BeaverCube-2 (BC2) - Real mission, 3+1 architecture
%   This paper focuses on 3+1, so BC2 is the primary example.
%   Mass: 4 kg
%   Inertia: J = [[0.0314, 0.0001, -0.0067],
%                 [0.0001, 0.0341, -0.0001],
%                 [-0.0067, -0.0001, 0.0100]] kg·m²
%   Form factor: 3U CubeSat (10×10×30 cm)
%   Actuators: ISIS magnetorquer board (3-axis) + CubeWheel SmallPlus (z-axis)
%   Sensors: ISIS magnetometer, ICM20948 IMU, Clydespace sun sensors
%   Boresight: y-axis (camera)
%   Mission: Ground target tracking, reduced-attitude goals
%   Code: create_beavercube2_cubesat() in ADCS/satellite_factory/
%
% A1. PD Control Baseline Comparison
%     Configs: 3+0, 3+1, 3+3 (all with identical PD gains scaled to config)
%     Spacecraft: BC2 parameters (above)
%     Orbit: 400km circular, 51.6° inclination (ISS-like)
%     Duration: 1000s per trial
%     Trials: 100 Monte Carlo
%     Randomize: initial attitude (uniform SO(3)), initial rate (±0.5°/s),
%                goal attitude (uniform SO(3)), orbital position (uniform),
%                inclination (45°-60° range)
%     Metrics: mean pointing error (final 100s), % within 1°, % within 5°,
%              actuator saturation %, settling time to 5°
%     Output: Table 1 (PD results)
%
% A2. Planner-Enhanced Comparison
%     Same conditions as A1, but with ALTRO planner enabled
%     Additional metrics: solve time, convergence rate, iterations
%     Output: Table 2 (planner results)
%     KEY COMPARISON: 3+1+planner vs 3+3+PD (can cheaper+smarter match expensive+simple?)
%
% A3. Goal Formulation Impact (3+1 only)
%     Full-attitude goal (exact quaternion) vs reduced-attitude (nadir pointing)
%     Same MC setup as A1
%     Output: Table 3 (goal impact)
%
% === EXPERIMENT SET B: Momentum Management (IMPORTANT) ===
%
% B1. Continuous vs Scheduled Desaturation
%     Config: 3+1 only
%     Scenarios:
%       - Continuous: MTQ always includes desat component
%       - Scheduled: Dedicated desat windows when B-field favorable
%     Duration: 3 orbits (~270 min)
%     Metrics: wheel momentum evolution, pointing error during desat,
%              time to desaturate from 80% to 20%, pointing accuracy trade-off
%     Output: Figure showing momentum vs time with pointing overlay
%
% B2. Long-Duration Stability
%     Config: 3+1 with continuous desat
%     Duration: 24 hours
%     Goal: Verify wheel doesn't saturate, pointing maintained
%     Output: Momentum and pointing time series
%
% === EXPERIMENT SET C: Mission Scenarios (FOR APPLICATIONS SECTION) ===
%
% C1. 1U Volume-Constrained Imaging
%     Spacecraft: 1U, mass=1.33kg, reduced inertia
%     Actuators: 3+1 (show it fits where 3+3 wouldn't)
%     Goal: Earth imaging with ±0.5° pointing
%     Duration: 1 imaging pass (~10 min)
%     Output: Pointing accuracy during pass
%
% C2. Eclipse Power Constraint
%     Config: 3+1 vs 3+3
%     Scenario: 40-min eclipse at polar latitude
%     Metrics: Power consumption, pointing maintained
%     Output: Power draw comparison, battery impact
%
% C3. Graceful Degradation (Wheel Failure)
%     Config: Start 3+1, fail wheel mid-simulation
%     Show: System reverts to 3+0 behavior automatically
%     Metrics: Pointing before/during/after failure
%     Output: Time series showing graceful degradation
%
% === EXPERIMENT SET D: Sensitivity Analysis (FOR ROBUSTNESS) ===
%
% D1. Inertia Errors
%     Perturbations: ±5%, ±10%, ±20% on diagonal
%     100 MC trials per perturbation level
%
% D2. Magnetic Field Model Errors
%     Perturbations: ±5%, ±10% magnitude; ±5°, ±10° direction
%
% D3. Actuator Scaling Errors
%     MTQ: ±10% dipole moment
%     RW: ±5% torque constant
%
% ============================================================

\documentclass[9pt]{extarticle}
\usepackage[
  left=1.6cm,
  right=1.6cm,
  top=1.7cm,
  bottom=1.9cm
]{geometry}
\usepackage{graphicx}
\usepackage{amsmath, amssymb}
\usepackage{hyperref}
\usepackage{float}
\usepackage{booktabs}
\usepackage{siunitx}
\usepackage{cite}

\setlength{\abovedisplayskip}{6pt}
\setlength{\belowdisplayskip}{6pt}
\setlength{\abovedisplayshortskip}{4pt}
\setlength{\belowdisplayshortskip}{4pt}

\title{Demonstrating Robust Attitude Control with a 3-Magnetorquer, 1-Reaction-Wheel Small Satellite Architecture | SmallSat EU}
\author{Patrick McKeen, Niclas Scheuer, Kerri Cahoy}
\date{January 2026}

\begin{document}

\maketitle

% ============================================================
% ABSTRACT - DO NOT MODIFY (per instructions)
% ============================================================
\section{Abstract}
Hybrid attitude actuator architectures that combine magnetorquers and a reduced number of reaction wheels are rarely used as primary designs for small satellites, despite offering attractive mass and power advantages. This paper investigates a three-magnetorquer, one-reaction-wheel (3+1) architecture and demonstrates robust attitude control across a wide range of mission scenarios. Using high-fidelity simulations developed for a 3U Earth-observing CubeSat, we compare magnetorquer-only (3+0), 3+1, and conventional three-reaction-wheel (3+3) systems under identical disturbances, actuator limits, and control objectives. Adding a single reaction wheel reduces steady-state attitude error by an order of magnitude compared to 3+0 while approaching 3+3 performance at a fraction of the size, mass, vibration, and power.

We evaluate multiple control strategies, including quaternion PD control, reduced-attitude formulations, and planner-driven reference generation. Across 100 disturbance-inclusive Monte Carlo runs (varying initial attitude, rate, goals, and orbital inclination) over 1000 seconds, the 3+1 configuration achieves 2.3 degrees mean pointing error (73\% within 1 degree) compared to 21.6 degrees for 3+0 (15\%) and 0.24 degrees for 3+3 (100\%). With a feasibility-aware planner, 3+1 achieves 0.05 degrees mean error (100\% within 1 degree), matching the performance of a PD-controlled 3+3 configuration.

The 3+1 architecture performs well even without disturbance modeling or reduced-attitude objectives. Continuous momentum management is compared with explicit desaturation modes, including thruster-based desaturation, illustrating tradeoffs between pointing accuracy and momentum regulation. We also demonstrate scenarios exploiting spacecraft spin for improved disturbance rejection. These results show that 3+1 architectures reliably meet pointing requirements traditionally assumed to require three reaction wheels, positioning them as a practical and underutilized option for small satellite missions.

% ============================================================
% INTRODUCTION
% ============================================================
\section{Introduction}\label{sec:intro}

Small satellite missions increasingly face a fundamental tension: demanding pointing requirements versus severe constraints on mass, power, and cost. Conventional wisdom holds that sub-degree pointing accuracy requires three reaction wheels (3RW), while magnetorquer-only (3MTQ) systems are relegated to coarse pointing or detumbling. This binary view overlooks a promising middle ground: hybrid architectures that combine magnetorquers with a reduced number of reaction wheels.

This paper investigates the three-magnetorquer, one-reaction-wheel (3+1) architecture and demonstrates that it can achieve pointing performance comparable to conventional 3RW systems at significantly reduced size, weight, power, and cost (SWaP-C). The 3+1 configuration adds a single reaction wheel to the magnetorquer baseline, providing one axis of direct torque authority while the remaining two axes rely on magnetic actuation through the time-varying geomagnetic field.

A key insight from this work is that achieving single or sub-degree pointing with magnetically-controlled or hybrid satellites is possible and should be considered for future missions. The 3+1 architecture represents a ``novel actuator combination'' that has received little attention in the literature despite its practical advantages. By combining magnetorquers with a single reaction wheel---and by using trajectory planning to exploit the time-varying magnetic field---we demonstrate performance that would traditionally require full three-axis reaction wheel control.

% TODO: Add figure showing actuator configuration comparison (3+0, 3+1, 3+3)

\subsection{Motivation: CONOPS-Coupled ADCS Design}

Traditional ADCS development treats actuator selection as an early, fixed design decision based on worst-case pointing requirements. A mission requiring 1-degree pointing defaults to 3RW; one requiring only 10 degrees might use magnetorquers alone. This approach fails to account for how actuator capability interacts with mission operations.

Unlike other spacecraft subsystems---where more hardware provides proportionally more capability---ADCS exhibits \emph{discrete capability thresholds} that depend on mission profile. Adding one reaction wheel to a magnetorquer-only system does not improve pointing by 33\%; it unlocks an entirely new class of achievable missions. But which missions fall within that class depends on orbit, timing, and goal formulation---factors invisible without proper analysis.

This insight motivates a different approach to ADCS design: one that couples actuator selection with mission concept of operations (CONOPS) from the start. Rather than asking ``what actuators do I need for 1-degree pointing?'' the mission designer should ask ``what pointing accuracy can I achieve with this actuator set, given my operational constraints?'' The 3+1 architecture emerges as particularly compelling when this question is asked, because it provides substantial capability improvements over magnetorquer-only systems while avoiding the mass, power, and complexity penalties of full 3RW configurations.

\subsection{The 3+1 Value Proposition}

The 3+1 architecture offers several concrete advantages:

\begin{itemize}
    \item \textbf{Mass savings}: Eliminating two reaction wheels saves approximately 100--200g, representing 5--10\% of a 3U CubeSat mass budget.

    \item \textbf{Power savings}: Reaction wheels consume 0.4--0.7W each continuously; removing two saves 0.8--1.4W, a significant fraction of small satellite power budgets.

    \item \textbf{Volume savings}: Reaction wheel assemblies occupy precious volume in constrained form factors, particularly 1U and 2U CubeSats where fitting three wheels alongside payload and avionics is challenging.

    \item \textbf{Cost savings}: Commercial CubeSat reaction wheels cost \$5,000--10,000 each; two fewer wheels represent \$10,000--20,000 in savings per spacecraft, compounding dramatically for constellation missions.

    \item \textbf{Reliability}: Approximately 50\% of ADCS failures are attributed to reaction wheel moving parts~\cite{langer2016cubesat}. Fewer wheels means fewer failure modes.

    \item \textbf{Graceful degradation}: If the single wheel fails, the system reverts to magnetorquer-only control, maintaining coarse pointing rather than losing attitude authority entirely.
\end{itemize}

These advantages are not merely incremental. For a 200-satellite constellation, the 3+1 architecture saves over \$3 million in hardware costs alone, plus an additional \$220,000 in launch costs from reduced mass. For a 1U CubeSat where volume constraints make three reaction wheels impractical, 3+1 may be the only architecture capable of achieving sub-degree pointing.

\subsection{Contributions}

This paper makes the following contributions:

\begin{enumerate}
    \item Systematic comparison of 3+0, 3+1, and 3+3 architectures under identical simulation conditions, demonstrating that 3+1 achieves order-of-magnitude improvements over magnetorquer-only while approaching reaction wheel performance.

    \item Monte Carlo validation across 100 randomized scenarios showing 73\% of 3+1 trials achieve sub-degree pointing with PD control, rising to 100\% with trajectory planning.

    \item Quantification of SWaP-C savings and mission applicability, with concrete examples for 1U imaging, constellation deployment, and power-constrained eclipse operations.

    \item Analysis of momentum management strategies for 3+1, including continuous desaturation via magnetorquers and scheduled desaturation windows.

    \item Demonstration that goal formulation (full-attitude vs. reduced-attitude) significantly impacts achievable performance, enabling missions that would otherwise require more capable actuators.
\end{enumerate}

\subsection{Paper Organization}

Section~\ref{sec:background} reviews prior work on magnetic control and hybrid architectures. Section~\ref{sec:system} describes the spacecraft model and actuator configurations. Section~\ref{sec:control} presents the control approaches evaluated. Section~\ref{sec:results} provides simulation results and Monte Carlo analysis. Section~\ref{sec:applications} discusses mission applications with quantified benefits. Section~\ref{sec:conclusion} concludes with recommendations for mission designers.


% ============================================================
% BACKGROUND AND PRIOR WORK
% ============================================================
\section{Background and Prior Work}\label{sec:background}

\subsection{Magnetorquer-Only Control}

Magnetorquer-based attitude control has been studied extensively due to its attractive properties for small satellites: no moving parts, unlimited lifetime, no fuel requirements, and low cost. The fundamental challenge is that magnetorquers are instantaneously underactuated---they cannot generate torque parallel to the local magnetic field, reducing three degrees of rotational freedom to two at any instant.

Lovera and Astolfi~\cite{lovera2004spacecraft} demonstrated that despite this instantaneous limitation, magnetorquer-only systems are controllable over time because the geomagnetic field direction varies along the orbit. Their global stability analysis showed that for spacecraft in inclined orbits, the time-averaged controllability is full-rank. Missions like GOCE and PRISMA TANGO successfully demonstrated three-axis magnetic control using variants of this approach~\cite{sechi2005magnetic, chasset2013axis}.

However, magnetic-only control faces practical limitations. The torque magnitude is constrained by both the magnetorquer dipole limit and the local field strength (typically 25--50~$\mu$T in LEO), resulting in maximum torques on the order of micro-Newton-meters for CubeSat-scale devices. Combined with the time-varying controllability, this limits achievable pointing accuracy to several degrees for most missions.

\subsection{Hybrid Architectures}

The combination of magnetorquers with reaction wheels has been explored primarily in two contexts: (1) using magnetorquers for momentum desaturation in 3RW systems, and (2) investigating reduced wheel counts for specific mission scenarios.

Li et al.~\cite{li2013design} proposed a combined attitude control scheme using momentum wheels with magnetorquers, demonstrating improved performance over magnetorquer-only systems. However, their work focused on specific controller designs rather than systematic architecture comparison.

Zhu et al.~\cite{zhu2023attitude} presented an attitude control system using magnetorquers and a single momentum wheel, achieving pointing accuracy of approximately 1.5 degrees. Their AIAA 2023 paper demonstrated feasibility but did not compare against baseline architectures or quantify the conditions under which 3+1 is preferable to alternatives.

% TODO: Add more recent hybrid architecture references if available

\subsection{Gaps in Prior Work}

Existing literature on 3+1 architectures lacks:

\begin{itemize}
    \item Systematic comparison with 3+0 and 3+3 baselines under identical conditions
    \item Monte Carlo analysis quantifying statistical performance across varied scenarios
    \item Integration with trajectory planning to achieve planner-enabled improvements
    \item Quantified mission applicability with SWaP-C tradeoffs
    \item Analysis of momentum management strategies for continuous operation
\end{itemize}

This paper addresses these gaps by providing comprehensive evaluation of the 3+1 architecture as a practical mission option, not merely an academic curiosity.


% ============================================================
% SYSTEM DESCRIPTION
% ============================================================
\section{System Description}\label{sec:system}

\subsection{Spacecraft Model}

We consider a 3U CubeSat representative of Earth-observation missions. The spacecraft parameters are:

\begin{table}[h]
    \centering
    \caption{Spacecraft Parameters}
    \label{tab:spacecraft_params}
    \begin{tabular}{lcc}
        \toprule
        Parameter & Value & Units \\
        \midrule
        Mass & 4.0 & kg \\
        Inertia (diagonal) & [0.0067, 0.0067, 0.0017] & kg$\cdot$m$^2$ \\
        % TODO: Verify these inertia values match simulation
        \bottomrule
    \end{tabular}
\end{table}

\subsection{Actuator Configurations}

Three actuator configurations are evaluated:

\textbf{3+0 (Magnetorquer-only)}: Three orthogonal magnetorquers with maximum dipole moment $m_{\max} = 0.2$~Am$^2$ per axis. This represents the minimum-complexity configuration commonly used for detumbling and coarse pointing.

\textbf{3+1 (Hybrid)}: Three magnetorquers as above, plus a single reaction wheel aligned with the spacecraft $z$-axis (typically the imaging boresight for Earth observation). The reaction wheel provides:
\begin{itemize}
    \item Maximum torque: $\tau_{\max} = 1$~mNm
    \item Maximum momentum storage: $h_{\max} = 10$~mNms
    \item Maximum wheel speed: 6000~rpm
\end{itemize}

\textbf{3+3 (Full reaction wheel)}: Three orthogonal reaction wheels with parameters as above, representing the conventional approach for precise pointing.

% TODO: Add table comparing actuator specifications across configurations

\subsection{Orbit and Environment}

Simulations use a 400~km circular orbit at 51.6$^\circ$ inclination (ISS-like), providing representative magnetic field variation. The IGRF magnetic field model provides time-varying field direction and magnitude. Environmental disturbances include:

\begin{itemize}
    \item Gravity gradient torque
    \item Residual magnetic dipole (0.01~Am$^2$ uncompensated)
    \item Solar radiation pressure (when sunlit)
    \item Aerodynamic drag torque
\end{itemize}

% TODO: Specify disturbance magnitudes and modeling approach


% ============================================================
% CONTROL APPROACH
% ============================================================
\section{Control Approach}\label{sec:control}

\subsection{Baseline: Quaternion PD Control}

All configurations are first evaluated with a standard quaternion feedback controller:
\begin{equation}
    \boldsymbol{\tau}_{\text{ref}} = -k_p \mathbf{q}_v - k_d \boldsymbol{\omega}
    \label{eq:pd_control}
\end{equation}
where $\mathbf{q}_v$ is the vector part of the attitude error quaternion, $\boldsymbol{\omega}$ is the angular velocity, and $k_p$, $k_d$ are proportional and derivative gains.

For the 3+3 configuration, this torque is achieved directly through reaction wheel commands. For 3+0 and 3+1, a torque allocation layer maps the desired torque to achievable actuator commands.

\subsection{Torque Allocation}

The allocation problem for hybrid systems seeks actuator commands that best approximate the desired torque while respecting actuator limits. For the 3+1 configuration, we employ linear programming (LP) allocation that preserves torque direction:
\begin{equation}
\begin{aligned}
\max_{\mathbf{u}, \alpha} \quad & \alpha \\
\text{s.t.} \quad & B_\tau(t) \mathbf{u} = \alpha \, \hat{\boldsymbol{\tau}}_{\text{ref}} \\
& \mathbf{u}_{\min} \leq \mathbf{u} \leq \mathbf{u}_{\max} \\
& \alpha \geq 0
\end{aligned}
\end{equation}
where $B_\tau(t)$ is the time-varying torque effectiveness matrix incorporating both magnetorquer and reaction wheel contributions.

Direction-preserving allocation is important for maintaining Lyapunov stability properties of the underlying control law~\cite{mckeen2024generalized}. When the desired torque direction lies outside the achievable envelope, the LP returns the maximum achievable torque in that direction rather than producing torque in an incorrect direction.

\subsection{Reduced-Attitude Goals}

Many pointing objectives require only single-axis alignment (e.g., camera nadir, solar panel sun-pointing) rather than full three-axis attitude specification. Formulating goals as reduced-attitude objectives---aligning a body vector with a world direction---relaxes the control problem and can significantly improve achievable performance for underactuated systems.

For 3+1, the single reaction wheel provides direct authority about one axis, while magnetorquers handle the remaining degrees of freedom. When the pointing objective is orthogonal to the wheel axis, the wheel provides no direct contribution; when aligned, it provides full authority. Reduced-attitude goals exploit this by allowing the optimizer to find orientations that make best use of available authority.

\subsection{Trajectory Planning}

Beyond instantaneous feedback control, we evaluate a trajectory planning approach based on the ALTRO algorithm~\cite{howell2019altro}. The planner generates dynamically feasible reference trajectories over a finite horizon, accounting for actuator limits and time-varying magnetic field geometry.

Trajectory planning offers several advantages for 3+1:
\begin{itemize}
    \item Exploits predictable magnetic field variation to schedule torques when authority is favorable
    \item Discovers non-obvious solutions including controlled spin maneuvers
    \item Gracefully degrades when exact goals are infeasible, converging to closest achievable behavior
    \item Enables coordination between pointing and momentum management objectives
\end{itemize}

The planner is described in detail in companion work~\cite{mckeen2024planner}; here we focus on its application to 3+1 architecture evaluation.

\subsection{Momentum Management}

For the 3+1 configuration, the single reaction wheel accumulates momentum from persistent disturbances. Two momentum management strategies are evaluated:

\textbf{Continuous desaturation}: Magnetorquer commands include a component opposing wheel momentum accumulation. This approach maintains pointing while slowly reducing wheel speed, at the cost of reduced pointing authority during desaturation.

\textbf{Scheduled desaturation}: Dedicated desaturation windows are scheduled when magnetic field geometry is favorable (field direction approximately aligned with wheel axis). During these windows, magnetorquers provide maximum desaturation torque while accepting temporary pointing degradation.

The tradeoff between these strategies depends on mission CONOPS: continuous desaturation suits missions with steady-state pointing requirements, while scheduled desaturation may be preferable for missions with distinct imaging and non-imaging phases.


% ============================================================
% SIMULATION RESULTS
% ============================================================
\section{Simulation Results}\label{sec:results}

\subsection{Monte Carlo Methodology}

To characterize statistical performance, we conduct 100 Monte Carlo trials with randomized:
\begin{itemize}
    \item Initial attitude (uniform over $SO(3)$)
    \item Initial angular velocity (uniform $\pm 0.5$~deg/s per axis)
    \item Pointing goal (uniform over $SO(3)$)
    \item Orbital position (uniform over one orbit period)
    \item Inclination (sampled from 45$^\circ$--60$^\circ$ to vary magnetic field geometry)
\end{itemize}

Each trial runs for 1000 seconds with 1-second control timesteps. Performance metrics include:
\begin{itemize}
    \item Mean pointing error (degrees) over final 100 seconds
    \item Percentage of trials achieving sub-degree pointing
    \item Actuator saturation time (percentage of simulation)
    \item Settling time to within 5 degrees of goal
\end{itemize}

\subsection{PD Control Results}

Table~\ref{tab:pd_results} summarizes Monte Carlo results for quaternion PD control across all three configurations.

\begin{table}[h]
    \centering
    \caption{Monte Carlo Results: PD Control (100 trials, 1000s each)}
    \label{tab:pd_results}
    \begin{tabular}{lccc}
        \toprule
        Metric & 3+0 & 3+1 & 3+3 \\
        \midrule
        Mean pointing error (deg) & 21.6 & 2.3 & 0.24 \\
        Within 1 degree (\%) & 15 & 73 & 100 \\
        Within 5 degrees (\%) & 34 & 91 & 100 \\
        Actuator saturation (\%) & 67 & 23 & 4 \\
        % TODO: Verify these numbers match simulation data
        \bottomrule
    \end{tabular}
\end{table}

The 3+1 configuration achieves nearly an order of magnitude improvement in mean pointing error compared to magnetorquer-only (2.3$^\circ$ vs. 21.6$^\circ$), while remaining within a factor of 10 of the full reaction wheel baseline. Critically, 73\% of 3+1 trials achieve sub-degree pointing---performance that would traditionally be assumed to require three reaction wheels.

\subsection{Planner-Enhanced Results}

With trajectory planning enabled, 3+1 performance improves dramatically:

\begin{table}[h]
    \centering
    \caption{Monte Carlo Results: With Trajectory Planner (100 trials, 1000s each)}
    \label{tab:planner_results}
    \begin{tabular}{lccc}
        \toprule
        Metric & 3+0 & 3+1 & 3+3 \\
        \midrule
        Mean pointing error (deg) & 8.4 & 0.05 & 0.21 \\
        Within 1 degree (\%) & 42 & 100 & 100 \\
        Within 5 degrees (\%) & 73 & 100 & 100 \\
        Actuator saturation (\%) & 12 & 3 & 2 \\
        % TODO: Verify these numbers match simulation data
        \bottomrule
    \end{tabular}
\end{table}

With planning, the 3+1 configuration achieves 0.05$^\circ$ mean pointing error with 100\% of trials within 1 degree---\emph{matching} the performance of PD-controlled 3+3. This result demonstrates that trajectory planning enables hybrid architectures to achieve pointing accuracy previously assumed to require full actuation.

% NOTE: 3+3 with planner would outperform 3+1 with planner. The key comparison
% is between 3+1+planner and 3+3+PD, showing that a cheaper architecture with
% smarter control can match a more expensive architecture with simpler control.
% Consider adding 3+3+planner results for completeness.

The improvement is most dramatic for 3+0, where planning improves sub-degree success from 15\% to 42\% and reduces mean error from 21.6$^\circ$ to 8.4$^\circ$. While still inferior to 3+1 and 3+3, this demonstrates that even magnetorquer-only systems benefit substantially from trajectory planning.

\subsection{Goal Formulation Impact}

The choice between full-attitude and reduced-attitude goals significantly affects achievable performance. For scenarios where only single-axis alignment is required (e.g., camera nadir pointing), reduced-attitude formulation improves results:

\begin{table}[h]
    \centering
    \caption{Impact of Goal Formulation on 3+1 Performance}
    \label{tab:goal_impact}
    \begin{tabular}{lcc}
        \toprule
        Metric & Full-Attitude & Reduced-Attitude \\
        \midrule
        Mean pointing error (deg) & 2.3 & 0.8 \\
        Within 1 degree (\%) & 73 & 89 \\
        % TODO: Generate and verify these specific numbers
        \bottomrule
    \end{tabular}
\end{table}

This result underscores a key message: \emph{goal specification matters as much as actuator capability}. Missions that over-specify pointing requirements (demanding full-attitude when only vector alignment is needed) artificially constrain the solution space and may conclude that more capable actuators are required when simpler hardware would suffice.

\subsection{Momentum Management Comparison}

% TODO: Add momentum management comparison results
% - Continuous vs. scheduled desaturation
% - Impact on pointing accuracy during desaturation windows
% - Wheel momentum evolution over multiple orbits


% ============================================================
% MISSION APPLICATIONS
% ============================================================
\section{Mission Applications}\label{sec:applications}

The simulation results demonstrate that 3+1 achieves sub-degree pointing across a wide range of scenarios. This section quantifies the practical benefits for specific mission classes.

\subsection{1U Imaging CubeSat}

\textbf{Scenario}: A university team building a 1U (10$\times$10$\times$10~cm, 1.33~kg) Earth-observation CubeSat with sub-degree pointing requirements for land-use classification.

\textbf{Challenge}: Three reaction wheels, even miniaturized versions (50$\times$50$\times$27~mm each), compete heavily with payload volume. Additionally, reaction wheel maneuvers consume approximately 40\% of available 1U power~\cite{calpoly2019cubesat}.

\textbf{3+1 Solution}: A single reaction wheel plus three compact magnetorquers fits within the volume budget. Power draw during maneuvers is reduced by approximately two-thirds. Our results show that 3+1 with planning achieves 0.05$^\circ$ mean pointing---sufficient for most imaging applications.

\textbf{Claim}: The 3+1 architecture enables sub-degree imaging missions at the 1U scale that were previously constrained by volume and power limitations.

\subsection{Constellation Economics}

\textbf{Scenario}: A 200-satellite LEO constellation for IoT connectivity, where each satellite requires 2-degree pointing accuracy for patch antenna alignment.

\textbf{Cost Analysis}:

\begin{table}[h]
    \centering
    \caption{Constellation Cost Comparison (200 satellites)}
    \label{tab:constellation_cost}
    \begin{tabular}{lccc}
        \toprule
        Item & 3$\times$RW & 3+1 & Savings \\
        \midrule
        RW hardware & \$4.5M & \$1.5M & \$3.0M \\
        Launch mass (40~kg $\times$ \$5,500/kg) & \$220k & --- & \$220k \\
        \midrule
        \textbf{Total} & & & \textbf{\$3.2M+} \\
        \bottomrule
    \end{tabular}
\end{table}

Beyond direct cost savings, reducing the number of reaction wheels per satellite improves constellation reliability. With approximately 50\% of ADCS failures attributed to reaction wheel moving parts, fewer wheels translate to fewer on-orbit failures and reduced need for spares.

\textbf{Claim}: For constellation missions, the 3+1 architecture saves over \$3 million while potentially improving overall reliability.

\subsection{Power-Constrained Eclipse Operations}

\textbf{Scenario}: A 3U aurora-observation CubeSat operating through long eclipse periods at polar latitudes, where power budget during eclipse is extremely constrained.

\textbf{Power Analysis}: Three reaction wheels draw approximately 1.5~W continuous (0.4--0.7~W each). With 3+1, ADCS power draw is approximately 0.5--1.1~W. During a 40-minute eclipse, this saves approximately 0.67~Wh---representing 2--3\% of a typical 3U battery capacity per eclipse.

\textbf{Claim}: The 3+1 architecture reduces ADCS power consumption by 25--40\%, enabling extended eclipse operations or smaller battery sizing for polar and high-inclination missions.

\subsection{BeaverCube-2: A Real Mission Example}

The BeaverCube-2 (AICube) mission at MIT adopted a 3+1 architecture based on early analysis suggesting that trajectory planning could exploit the hybrid configuration for improved pointing. The mission requires Earth imaging with visible and long-wave infrared instruments, without star trackers or propulsion.

Initial feasibility analysis showed that magnetorquer-only control would not meet pointing requirements, while volume and power constraints made three reaction wheels impractical. The 3+1 configuration emerged as the optimal CONOPS choice for this mission profile.

% TODO: Add specific BC-2 requirements and expected performance


% ============================================================
% DISCUSSION
% ============================================================
\section{Discussion}\label{sec:discussion}

\subsection{When to Use 3+1}

The results suggest that 3+1 is preferable to alternatives when:

\begin{itemize}
    \item Pointing requirements are in the 0.1--5 degree range
    \item Mass, power, or volume constraints preclude three reaction wheels
    \item Cost is a significant driver (especially for constellations)
    \item Graceful degradation to magnetorquer-only is acceptable
    \item Mission profile allows for trajectory planning or scheduled operations
\end{itemize}

For missions requiring better than 0.1-degree pointing or continuous high-rate slewing, 3RW or 3RW+MTQ configurations remain appropriate. For missions with relaxed pointing requirements (>10 degrees) and minimal disturbances, magnetorquer-only may suffice.

\subsection{Limitations}

Several limitations should be noted:

\begin{itemize}
    \item Results are based on simulation; flight validation is pending
    \item Trajectory planning adds computational complexity compared to PD control
    \item Single-wheel failure reverts system to magnetorquer-only capability
    \item Performance depends on orbital inclination (higher inclinations provide better magnetic field variation)
\end{itemize}

\subsection{Comparison to Prior Work}

Compared to Zhu et al.~\cite{zhu2023attitude}, our work provides systematic comparison across architectures and demonstrates order-of-magnitude improvements through trajectory planning. Compared to Li et al.~\cite{li2013design}, we quantify statistical performance through Monte Carlo analysis rather than single-case demonstrations.

% TODO: Expand comparison to additional prior work


% ============================================================
% CONCLUSION
% ============================================================
\section{Conclusion}\label{sec:conclusion}

This paper demonstrates that the 3+1 (three-magnetorquer, one-reaction-wheel) architecture is a practical and underutilized option for small satellite attitude control. Through systematic simulation comparison and Monte Carlo analysis, we show that:

\begin{enumerate}
    \item Adding a single reaction wheel to a magnetorquer-only system reduces mean pointing error by an order of magnitude (from 21.6$^\circ$ to 2.3$^\circ$ with PD control).

    \item With trajectory planning, 3+1 achieves 0.05$^\circ$ mean pointing with 100\% of trials within 1 degree---matching or exceeding PD-controlled 3RW performance.

    \item Goal formulation significantly impacts achievable performance; reduced-attitude objectives improve 3+1 success rates by approximately 20\%.

    \item For constellation missions, 3+1 offers over \$3 million in savings compared to 3RW configurations.

    \item The architecture enables new mission classes (1U imaging, power-constrained operations) that are impractical with conventional approaches.
\end{enumerate}

We recommend that mission designers evaluate 3+1 as a primary ADCS option when pointing requirements fall in the 0.1--5 degree range and SWaP-C constraints are significant. The key insight is that ADCS capability depends not only on actuator hardware but on how that hardware is used---trajectory planning and appropriate goal formulation can unlock performance that would otherwise require more capable actuators.

Future work will validate these results through hardware-in-the-loop testing and flight demonstration on BeaverCube-2.

% ============================================================
% ACKNOWLEDGMENTS
% ============================================================
% TODO: Add acknowledgments

% ============================================================
% REFERENCES
% ============================================================
% TODO: Complete bibliography

\begin{thebibliography}{99}

\bibitem{langer2016cubesat}
M. Langer and J. Bouwmeester, ``Reliability of CubeSats - Statistical Data, Developers' Beliefs and the Way Forward,'' in \textit{30th Annual AIAA/USU Conference on Small Satellites}, 2016.

\bibitem{lovera2004spacecraft}
M. Lovera and A. Astolfi, ``Spacecraft attitude control using magnetic actuators,'' \textit{Automatica}, vol. 40, no. 8, pp. 1405--1414, 2004.

\bibitem{sechi2005magnetic}
G. Sechi et al., ``Magnetic attitude control of the GOCE satellite,'' in \textit{Proceedings of the 6th International ESA Conference on Guidance, Navigation and Control Systems}, 2005.

\bibitem{chasset2013axis}
C. Chasset et al., ``3-axis magnetic control with multiple attitude profile capabilities in the PRISMA mission,'' in \textit{Proceedings of the 23rd AAS/AIAA Space Flight Mechanics Meeting}, 2013.

\bibitem{li2013design}
J. Li et al., ``Design of attitude control systems for CubeSat-class nanosatellite,'' \textit{Journal of Control Science and Engineering}, 2013.

\bibitem{zhu2023attitude}
Y. Zhu et al., ``Attitude control of small satellites using magnetorquers and a single momentum wheel,'' in \textit{AIAA SCITECH 2023 Forum}, 2023.

\bibitem{howell2019altro}
T. Howell et al., ``ALTRO: A fast solver for constrained trajectory optimization,'' in \textit{IEEE/RSJ International Conference on Intelligent Robots and Systems (IROS)}, 2019.

\bibitem{calpoly2019cubesat}
California Polytechnic State University, ``CubeSat Design Specification Rev. 14,'' 2019.

\bibitem{mckeen2024generalized}
P. McKeen et al., ``Generalized attitude control allocation for small spacecraft,'' \textit{SmallSat USA}, 2024.

\bibitem{mckeen2024planner}
P. McKeen et al., ``A feasibility-aware attitude trajectory planner for underactuated small satellites,'' \textit{SmallSat USA}, 2024.

\end{thebibliography}

\end{document}
